% ---------------------------------------------------------------------------
% Corpus, task definition and deployed system.
% \input{sections/03_corpus_system} — expects natbib (\citep) and, for
% Figure 1, TikZ with libraries: arrows.meta, positioning.
% Counts here are the final production figures; accuracy lives in §Results.
% ---------------------------------------------------------------------------

\section{Corpus and Deployed System}
\label{sec:corpus}

\subsection{Corpus, task and output contract}
\label{ssec:corpus-task}

The system reads UNDRR's internal document lake, the text layer behind
PreventionWeb. Its contents are submitted rather than collected: national
disaster risk reduction strategies and Sendai Framework reporting from member
states, guidance material from UN agencies, project reports and evaluations
from NGOs, post-disaster needs assessments, training documentation. Events
appear inside them as case studies, annexes and background, not as the
document's subject. UNDRR holds no redistribution rights over the knowledge
base, so the corpus cannot be released with the paper; we report its
provenance, composition and access constraints as a data statement, and
document the derived record table separately
\citep{bender2018DataStatements,gebru2021Datasheets}.

The unit of processing is one document. The hazard is often named in a
heading, the date in a preceding clause, the affected districts two paragraphs
later; treating the document rather than the sentence as the extraction
context is the established response to that dispersion
\citep{li2021WikiEvents,ebner2020RAMS,tong2022DocEE}. Each document is
truncated to its first 12{,}000 characters, a cost-driven cutoff we have not
validated against a full-length re-run (\S\ref{sec:limitations}).

The output contract is a list of event records, possibly empty. Each carries a
hazard type from a controlled vocabulary, one or more event dates, one or more
affected-location strings with a parallel self-tagged administrative level, a
stated confidence, up to five impact statements, up to three short evidence
snippets from the source, and an explicit \texttt{missing\_information} list.
Hazard labels are constrained to the PreventionWeb vocabulary --- 19 main
labels plus five collection labels --- which operationalises the UNDRR--ISC
Hazard Information Profiles, a multi-hazard classification developed through
global expert consultation \citep{murray2021HazardProfiles}. The prompt
supplies the mapping rules explicitly: \emph{typhoon} and \emph{storm surge}
map to \textsc{Cyclone, Hurricane and Typhoon}, dam and transport failures to
\textsc{Technological hazard}. Compound events join exact labels with a pipe
(\texttt{Earthquake | Tsunami}); a real event with no exact label is dropped
rather than approximated.

A record is \emph{verifiable} only if it has a hazard label, at least one
event date, at least one location passing the location rules, and stated
confidence of high or medium. Three things are excluded by construction:
general risk statements, future projections and scenarios, and preparedness
activity with no actual event behind it. Publication and workshop dates are
likewise excluded from \texttt{event\_dates} --- the densest dates in these
documents, and the easiest to mistake for the event.

\subsection{Extraction and geocoding pipeline}
\label{ssec:pipeline}

Documents are pulled by one SQL query over two Databricks tables, joining
per-asset text to chunk-level metadata and filtering empty text. Preprocessing
is truncation and nothing else: no sentence splitting, no hazard keyword
filter, no genre pre-classification. Every retrieved document reaches the
model.

% ---------------------------------------------------------------------------
% Figure 1 — end-to-end flow, single column, compact snaking layout.
% % ---------------------------------------------------------------------------
% Figure 1 — end-to-end flow, single column, compact snaking layout.
% % ---------------------------------------------------------------------------
% Figure 1 — end-to-end flow, single column, compact snaking layout.
% \input{figures/fig_pipeline} from sections/03_corpus_system.tex.
% Requires: \usepackage{tikz} with \usetikzlibrary{arrows.meta,positioning}
% Colours are defined here so the file is self-contained.
% ---------------------------------------------------------------------------

\definecolor{pipeInk}{HTML}{1F3A5F}
\definecolor{pipeFill}{HTML}{EEF2F7}
\definecolor{pipeOut}{HTML}{2A6F6B}

\begin{figure}[t]
\centering
\begin{tikzpicture}[
    every node/.style={font=\scriptsize},
    stage/.style={
        rectangle, rounded corners=2pt,
        draw=pipeInk, line width=0.5pt, fill=pipeFill,
        text width=17mm, minimum height=9mm, align=center, inner sep=1.5pt
    },
    output/.style={
        rectangle, rounded corners=2pt,
        draw=pipeOut, line width=0.7pt, fill=pipeOut!10,
        text width=17mm, minimum height=9mm, align=center, inner sep=1.5pt
    },
    arr/.style={-{Stealth[length=1.6mm]}, line width=0.5pt, draw=pipeInk}
]

\node[stage] (doc) {Document\\{\tiny 12k chars}};
\node[stage, right=4mm of doc] (extract)
  {Schema-con\-strained extraction};
\node[stage, right=4mm of extract] (rewrite)
  {Location rewriting};
\node[stage, below=6mm of extract] (geo)
  {Geocoding\\{\tiny Nominatim/OSM}};
\node[output, below=6mm of doc] (rec)
  {Event record\\{\tiny + evidence}};

\draw[arr] (doc)     -- (extract);
\draw[arr] (extract) -- (rewrite);
\draw[arr] (rewrite.south) |- ([yshift=0mm]geo.east);
\draw[arr] (geo) -- (rec);

\node[font=\tiny\itshape, text=pipeInk!75, above=0.5mm of extract]
  {controlled hazard vocabulary};
\node[font=\tiny\itshape, text=pipeInk!75, below=0.5mm of geo]
  {polygon, EPSG:4326};

\end{tikzpicture}
\caption{The deployed pipeline. One document produces zero or more event
records; each retains the evidence snippets, source identifiers and resolution
trace that justify it.}
\label{fig:pipeline}
\end{figure}
 from sections/03_corpus_system.tex.
% Requires: \usepackage{tikz} with \usetikzlibrary{arrows.meta,positioning}
% Colours are defined here so the file is self-contained.
% ---------------------------------------------------------------------------

\definecolor{pipeInk}{HTML}{1F3A5F}
\definecolor{pipeFill}{HTML}{EEF2F7}
\definecolor{pipeOut}{HTML}{2A6F6B}

\begin{figure}[t]
\centering
\begin{tikzpicture}[
    every node/.style={font=\scriptsize},
    stage/.style={
        rectangle, rounded corners=2pt,
        draw=pipeInk, line width=0.5pt, fill=pipeFill,
        text width=17mm, minimum height=9mm, align=center, inner sep=1.5pt
    },
    output/.style={
        rectangle, rounded corners=2pt,
        draw=pipeOut, line width=0.7pt, fill=pipeOut!10,
        text width=17mm, minimum height=9mm, align=center, inner sep=1.5pt
    },
    arr/.style={-{Stealth[length=1.6mm]}, line width=0.5pt, draw=pipeInk}
]

\node[stage] (doc) {Document\\{\tiny 12k chars}};
\node[stage, right=4mm of doc] (extract)
  {Schema-con\-strained extraction};
\node[stage, right=4mm of extract] (rewrite)
  {Location rewriting};
\node[stage, below=6mm of extract] (geo)
  {Geocoding\\{\tiny Nominatim/OSM}};
\node[output, below=6mm of doc] (rec)
  {Event record\\{\tiny + evidence}};

\draw[arr] (doc)     -- (extract);
\draw[arr] (extract) -- (rewrite);
\draw[arr] (rewrite.south) |- ([yshift=0mm]geo.east);
\draw[arr] (geo) -- (rec);

\node[font=\tiny\itshape, text=pipeInk!75, above=0.5mm of extract]
  {controlled hazard vocabulary};
\node[font=\tiny\itshape, text=pipeInk!75, below=0.5mm of geo]
  {polygon, EPSG:4326};

\end{tikzpicture}
\caption{The deployed pipeline. One document produces zero or more event
records; each retains the evidence snippets, source identifiers and resolution
trace that justify it.}
\label{fig:pipeline}
\end{figure}
 from sections/03_corpus_system.tex.
% Requires: \usepackage{tikz} with \usetikzlibrary{arrows.meta,positioning}
% Colours are defined here so the file is self-contained.
% ---------------------------------------------------------------------------

\definecolor{pipeInk}{HTML}{1F3A5F}
\definecolor{pipeFill}{HTML}{EEF2F7}
\definecolor{pipeOut}{HTML}{2A6F6B}

\begin{figure}[t]
\centering
\begin{tikzpicture}[
    every node/.style={font=\scriptsize},
    stage/.style={
        rectangle, rounded corners=2pt,
        draw=pipeInk, line width=0.5pt, fill=pipeFill,
        text width=17mm, minimum height=9mm, align=center, inner sep=1.5pt
    },
    output/.style={
        rectangle, rounded corners=2pt,
        draw=pipeOut, line width=0.7pt, fill=pipeOut!10,
        text width=17mm, minimum height=9mm, align=center, inner sep=1.5pt
    },
    arr/.style={-{Stealth[length=1.6mm]}, line width=0.5pt, draw=pipeInk}
]

\node[stage] (doc) {Document\\{\tiny 12k chars}};
\node[stage, right=4mm of doc] (extract)
  {Schema-con\-strained extraction};
\node[stage, right=4mm of extract] (rewrite)
  {Location rewriting};
\node[stage, below=6mm of extract] (geo)
  {Geocoding\\{\tiny Nominatim/OSM}};
\node[output, below=6mm of doc] (rec)
  {Event record\\{\tiny + evidence}};

\draw[arr] (doc)     -- (extract);
\draw[arr] (extract) -- (rewrite);
\draw[arr] (rewrite.south) |- ([yshift=0mm]geo.east);
\draw[arr] (geo) -- (rec);

\node[font=\tiny\itshape, text=pipeInk!75, above=0.5mm of extract]
  {controlled hazard vocabulary};
\node[font=\tiny\itshape, text=pipeInk!75, below=0.5mm of geo]
  {polygon, EPSG:4326};

\end{tikzpicture}
\caption{The deployed pipeline. One document produces zero or more event
records; each retains the evidence snippets, source identifiers and resolution
trace that justify it.}
\label{fig:pipeline}
\end{figure}


Five prompts were compared during development, varying location strictness,
case-study recall, an explicit verification step and schema compactness
(\S\ref{sec:prompts}); one was frozen for production and every number in this
paper comes from it. Extraction is one call per document to an
instruction-tuned model, with no fine-tuning and no in-domain event labels to
fine-tune on; the prompt carries the label semantics and decision rules a
trained model would otherwise absorb from data
\citep{sainz2024GoLLIE,jiao2023InstructIE}. The call is constrained to a
Pydantic model exposed to the API as a JSON schema, following work that uses
formal output constraints and typed class definitions for structured tasks
\citep{geng2023GrammarConstrained,wang2023Code4Struct}; the constraint is
warranted because unconstrained models still fail schema validity at
non-trivial rates \citep{lu2025SchemaRL}. It buys parseability, not truth, and
may itself shift model behaviour \citep{schall2025HiddenCost} --- which is why
\S\ref{sec:validation} judges records against sources rather than against the
schema. A call that fails after retry is written as an error row rather than
dropped, so a resume re-queues it; post-parse validators coerce the parallel
admin-level list to the allowed tags and drop non-verifiable, low-confidence
and duplicate records.

Locations are resolved behind a swappable interface, \texttt{resolve(text,
country\_hint)}, because geoparser choice is a known source of incomparability
and production will not keep this one
\citep{gritta2018Missing,halterman2023Mordecai3}. The implementation queries
Nominatim over OpenStreetMap, throttled and cached. A raw mention rarely
matches OSM directly, so each is walked down a rewrite ladder: the original
string; a possessive rewrite splitting owner from feature and converting the
owner to a country code (\emph{Thailand's Andaman coast} $\rightarrow$
\emph{Andaman coast}, \texttt{th}); facility-noun rewrites (\emph{Beirut
seaport} $\rightarrow$ \emph{Port of Beirut}); a leading-qualifier strip
(\emph{northern}, \emph{areas around}); and last, the owner country alone.
Where the document names exactly one country it is passed as a hint;
multi-country reports pass none. This country-first ordering follows resolvers
that fix larger administrative units before using them as context for smaller
ones \citep{zhang2023GeoNorm}, and targets the failure mode that dominates our
error analysis: one string denoting places in different countries
\citep{gritta2018WhichMelbourne}.

The resolver returns a Polygon or MultiPolygon in EPSG:4326, falling back to a
bounding box or a population-scaled point buffer where OSM has no polygon. The
country-level last resort is flagged and its confidence capped, so a coarse
geometry stays visible as coarse rather than passing as a resolved district.
Every emitted row keeps its evidence snippets, source identifier and URL,
prompt and schema version, the ladder step and country hint behind the query,
the returned administrative level, and the geometry's provenance. That trace
is the audit artifact: a downstream user can contest the path from a sentence
in a national report to a polygon on a map \citep{raji2020InternalAuditing}.
